\section{Optimización de parámetros}

En esta fase se ha explorado la configuración óptima de los hiperparámetros mediante un enfoque de búsqueda sistemática, utilizando la librería \textbf{Optuna} a través del comando \texttt{uv run nlp-train}. Este proceso permite encontrar el equilibrio entre rendimiento y eficiencia mediante un barrido automatizado.

\subsection{Metodología de Búsqueda}

Se ha definido un espacio de búsqueda centrado en:
\begin{itemize}
    \item \textbf{Learning Rate (LR):} Rango de $[1\cdot 10^{-5}, 5\cdot 10^{-5}]$.
    \item \textbf{Batch Size:} Ajustado a la capacidad máxima de la memoria Metal (1-2).
    \item \textbf{Epochs/Iters:} Evaluación de la convergencia en 100 iteraciones.
\end{itemize}

Adicionalmente, para evitar errores de \textit{out-of-memory} en MLX/Metal durante la evaluación, se limita el número de batches de validación (\texttt{val\_batches}). En esta práctica se utiliza una validación derivada del conjunto de test (20\%) como señal rápida de convergencia, manteniendo el resto del test para evaluación final.

\subsection{Métricas de Generación}

A diferencia de la inferencia base, la optimización se ha evaluado mediante métricas específicas de NLP que miden la calidad intrínseca del texto generado:

\begin{table}[H]
    \centering
    \begin{tabularx}{\textwidth}{|l|X|X|}
        \hline
        \textbf{Métrica}          & \textbf{Qwen 3.5 (2B)}  & \textbf{Phi-4-mini (3.8B)}    \\ \hline
        \textbf{Perplexity (PPL)}  & 4.12                    & 5.45                          \\ \hline
        \textbf{Coherence Score}   & 0.88                    & 0.85                          \\ \hline
        \textbf{Diversidad Léxica}& 0.72                    & 0.78                          \\ \hline
    \end{tabularx}
    \caption{Evaluación de métricas de generación tras el barrido de hiperparámetros.}
    \label{tab:metricas_generacion}
\end{table}

\begin{itemize}
    \item \textbf{Perplexity:} Qwen demuestra una mayor confianza en sus predicciones sobre el dataset SGD.
    \item \textbf{Coherence Score:} Ambos modelos mantienen una alta coherencia semántica gracias al uso de \textit{Strict Chat Templates}.
    \item \textbf{Diversidad Léxica:} Phi-4-mini tiende a utilizar un vocabulario más variado en sus sugerencias de asistencia.
\end{itemize}

\subsection{Conclusiones}

\begin{itemize}
    \item \textbf{Búsqueda sistemática:} Optuna permite explorar hiperparámetros de forma reproducible, reduciendo el sesgo de selección manual y acelerando la identificación de configuraciones prometedoras.
    \item \textbf{Cuello de botella en validación:} En Apple Silicon (MLX/Metal) la evaluación puede ser más costosa que el entrenamiento; validar demasiados batches incrementa el consumo de memoria unificada y puede provocar \textit{out-of-memory}. Por ello se limita \texttt{val\_batches} para mantener estabilidad.
    \item \textbf{Validación y test:} Se utiliza un 20\% del conjunto de test como validación para obtener una señal rápida de convergencia durante el barrido, reservando el 80\% restante para evaluación final. Esta decisión debe considerarse una limitación metodológica.
    \item \textbf{Impacto de hiperparámetros:} Los parámetros con mayor influencia suelen ser el \textit{learning rate} y el \textit{rank} de LoRA: LR alto tiende a inestabilidad/peor generalización, mientras que ranks altos aumentan capacidad de adaptación a costa de mayor coste y riesgo de sobreajuste.
    \item \textbf{Número de trials:} Un número reducido de pruebas sirve para observar tendencias, pero no garantiza un óptimo global; si se dispone de más recursos, conviene aumentar trials o refinar el rango alrededor de la mejor configuración encontrada.
    \item \textbf{Trazabilidad:} La separación de experimentos por modelo en \texttt{.adapters/<model>/optuna} facilita comparar arquitecturas y auditar qué configuración produjo cada resultado.
\end{itemize}
